\documentclass[10pt]{article}
\usepackage[margin=20mm]{geometry}
\usepackage{amsmath,amssymb,mathtools,fontspec,unicode-math,xurl,hyperref}
\IfFontExistsTF{DejaVu Serif}{\setmainfont{DejaVu Serif}}{\setmainfont{Latin Modern Roman}}
\setmathfont{latinmodern-math.otf}
\setlength{\parindent}{0pt}\setlength{\parskip}{6pt}\sloppy
\title{Scientific review of M25-0021:
a linearization coefficient in rarefaction-wave energy estimates}
\author{}
\date{25 October 2026}
\begin{document}
\maketitle
Abstract: This claim-focused initial review checks the local Euler linearization, its use in the energy proof, and the refined Gronwall lemma in the manifest-pinned author manuscript. A factor-of-two error in the linearized wave equation is established both analytically and by an exact smooth Euler solution. It propagates to the displayed Riemann-invariant wave system used in the core energy argument. The supported judgment is Major for this invalid core derivation, not a counterexample to the main a priori theorem. A local correction is derived; the complete corrected bootstrap proof has not been verified. No issue was found in the inspected Gronwall argument.\par
\section*{Identity, version, and scope}
The reviewed paper is Tian-Wen Luo and Pin Yu, On the stability of multi-dimensional rarefaction waves I: the energy estimates, arXiv:2302.09714v2, revised 19 September 2024. The supplied PDF has 100 pages; page references below are physical, one-based PDF pages (matching the printed numbers at the cited locations). The journal identifier is DOI 10.4007/annals.2025.202.2.2. This review binds to the supplied author manuscript, not to an independently compared journal final version. Their equivalence has not been established.\par
All 30 manifest-listed input files matched their recorded SHA-256 hashes. The PDF hash is\par
> 78ce7a2d283e922d47e1826882407b530f2986ce0843da306bcd016f508a6e78.\par
The principal TeX source is `I-Rarefaction\_waves\_final-version.tex`, with 4970 lines and SHA-256\par
> b0664c3470ca73d7aa77c686d6a98d844d48de95b75a7bbc4e978d86fff93925.\par
The actual text-reading scope was PDF pp. 1–8, 23–38, 50–55, 65–67, and 89–91. Source line ranges 790–820, 1058–1075, and 1385–1420 were read. Actual page images inspected were pp. 1, 23, 28, 33, 34, 37, 53–55, 65–66, and 89–90. Extraction of the remaining pages does not constitute reading or visual inspection. This is not a certification of the full 100-page proof.\par
\section*{Definitions, quantifiers, and the main theorem}
The paper studies two spatial dimensions, with the second coordinate periodic, isentropic pressure \(p=k_0\rho^\gamma\), \(k_0>0\), and \(1<\gamma<3\). In the irrotational potential formulation, \[v=-\nabla\phi,\qquad
h=\phi_t-\tfrac12|\nabla\phi|^2=\frac{c^2}{\gamma-1},
\qquad c^2=p'(\rho).\] Here \(\psi_\mu=\partial_\mu\phi\), in particular \(\psi_i=-v^i\) for \(i=1,2\). The acoustic metric and scalar contraction are \[g=-c^2dt^2+(dx^1-v^1dt)^2+(dx^2-v^2dt)^2,\qquad
Q_g(f,k)=g^{\mu\nu}\partial_\mu f\,\partial_\nu k.\] Equation (2.16), p. 28, defines \[w=\frac{c}{\gamma-1}+\frac{\psi_1}{2},\qquad
\underline{w}=\frac{c}{\gamma-1}-\frac{\psi_1}{2}.\] The main theorem on p. 37 is an a priori estimate under the specific rarefaction data conditions \((I_0),(I_2),(I_\infty)\) and the preceding irrotational setting. It requires \(N_{\rm top}\geq9\) and gives an \(\varepsilon_0\) and energy constants independent of \(0<\delta<1/2\). Its conclusion includes \(t^2\varepsilon^2\) energy control. These special data assumptions are much stronger than the local smooth Euler assumptions under which the identities in Section 2 are derived. The example below tests those local identities; it is not asserted to meet the global rarefaction-data ansatz.\par
\section*{Problem 1: an incorrect coefficient in the local wave system}
\subsection*{Source claim and precise location}
Equation reference: claimed, claimedw. On p. 23, Section 2 (TeX lines 793–811), the manuscript sets \(\widetilde g=\Omega g\) with \(\Omega=\rho/c\), asserts \(\Box_{\widetilde g}\psi=0\) for a variation of the potential, and displays \[
\Box_g\psi=-\tfrac12Q_g(\log(\rho/c),\psi)
\qquad\text{[paper (2.2)].}\] On pp. 33–34, paper (3.7) gives, in particular, \[
\Box_gw=-c^{-1}\left[
\frac{\gamma+1}{4}Q_g(\underline{w},\underline{w})+
\frac{\gamma-3}{4}Q_g(\underline{w},w)+
Q_g(w,w)+\frac12Q_g(\psi_2,\psi_2)\right].\] Its \(\psi_2\) equation has coefficient \((3-\gamma)/(4c)\) multiplying \(-Q_g(w,\psi_2)-Q_g(\underline{w},\psi_2)\). These are actual displayed formulas, confirmed against the rendered PDF and TeX lines 1403–1411, not extraction artifacts.\par
\subsection*{Complete analytic check}
Equation reference: divergence, correct. For the Lagrangian \(p(h)\), differentiation with respect to \(h\) gives \(p_h=\rho\) and \(p_{hh}=\rho/c^2\). The Hessian with respect to the three components of \(d\phi\) is \[L_{\phi_\mu\phi_\nu}=-\rho g^{\mu\nu}.\] For example, its \(00\) entry is \(\rho/c^2\), its \(0i\) entry is \(\rho v^i/c^2\), and its \(ij\) entry is \(\rho v^iv^j/c^2-\rho\delta_{ij}\). Thus the linearized Euler–Lagrange equation is \[
\partial_\mu(\rho g^{\mu\nu}\partial_\nu\psi)=0.\] The stated acoustic metric has \(\det g=-c^2\), so \[\Box_g\psi=c^{-1}\partial_\mu(cg^{\mu\nu}\partial_\nu\psi).\] Dividing [divergence] (\#divergence) by \(\rho\) proves the corrected local identity \[
\boxed{\ \Box_g\psi=-Q_g(\log(\rho/c),\psi).\ }\] The coefficient is one, not one half.\par
The same discrepancy follows directly from conformal scaling in three space-time dimensions. For a positive \(F\), \[\sqrt{|\det(Fg)|}\,(Fg)^{\mu\nu}
=F^{1/2}c\,g^{\mu\nu}.\] To obtain the density \(\rho g^{\mu\nu}\) in [divergence] (\#divergence), one must take \(F=(\rho/c)^2\). Consequently, retaining the manuscript’s \(\Omega=\rho/c\) requires \(\widetilde g=\Omega^2g\). Equivalently one could redefine \(\Omega=(\rho/c)^2\) while retaining \(\widetilde g=\Omega g\), but all subsequent occurrences of \(d\log\Omega\) would then need correction.\par
Equation reference: correctw. For completeness the general corrected Riemann-invariant system is also derivable without guessing coefficients. Set \(b=\gamma-1\). Since \(\log(\rho/c)=((3-\gamma)/b)\log c+\text{constant}\), [correct] (\#correct) gives \[\Box_g\psi_a=-\frac{3-\gamma}{bc}Q_g(c,\psi_a),\qquad a=0,1,2.\] Applying the product rule to \(h=\psi_0-(\psi_1^2+\psi_2^2)/2=c^2/b\) yields \[\Box_gc=-\frac{2}{bc}Q_g(c,c)
-\frac{b}{2c}\{Q_g(\psi_1,\psi_1)+Q_g(\psi_2,\psi_2)\}.\] Substituting \(c=b(w+\underline{w})/2\) and \(\psi_1=w-\underline{w}\) then gives \[\begin{aligned}
\Box_gw&=-c^{-1}\left[
 \frac{7-\gamma}{4}Q_g(w,w)+\frac{\gamma+1}{4}Q_g(\underline{w},\underline{w})
 +\frac12Q_g(\psi_2,\psi_2)\right],\\
\Box_g\underline{w}&=-c^{-1}\left[
 \frac{\gamma+1}{4}Q_g(w,w)+\frac{7-\gamma}{4}Q_g(\underline{w},\underline{w})
 +\frac12Q_g(\psi_2,\psi_2)\right],\\
\Box_g\psi_2&=-\frac{3-\gamma}{2c}
 \{Q_g(w,\psi_2)+Q_g(\underline{w},\psi_2)\}.\end{aligned}\] These are verified local corrections, not a verified replacement proof of the main theorem.\par
\subsection*{An exact admissible test of the local assertion}
Equation reference: example. Let \(0<t<1\), \(x=x^1\in\mathbb R\), \(y=x^2\in\mathbb R/(2\pi\mathbb Z)\), and \(s=1+t\). Choose \[
\gamma=2,\quad k_0=\tfrac12,\quad
\rho=s^{-1},\quad p=\tfrac12s^{-2},\quad c=s^{-1/2},\quad
v=(x/s,0),\quad \phi=-\frac{x^2}{2s}+\log s .\] All fields are smooth on this domain, \(\rho,c>0\), \(1<\gamma<3\), \(k_0>0\), and the fields are periodic in \(y\). The flow is irrotational and \(v=-\nabla\phi\). Writing \(B=\partial_t+v\cdot\nabla\), one checks every local Euler condition: \[B\rho=-s^{-2},\quad \rho\,\operatorname{div}v=s^{-2},\qquad
Bv=(0,0),\quad \rho^{-1}\nabla p=(0,0).\] Furthermore \[\phi_t-\tfrac12|\nabla\phi|^2=s^{-1}=c^2/(\gamma-1).\] The nondivergence potential equation \(g^{\mu\nu}\phi_{\mu\nu}=0\) also vanishes identically. The family \(\phi_\lambda(t,x,y)=\phi(t+\lambda,x,y)\) is an exact local family of such solutions for sufficiently small \(\lambda\). Its variation is \(\psi_0=\phi_t\), so it falls under the manuscript’s stated linearization.\par
Direct calculation from the coordinate divergence formula for \(\Box_g\), independently of either claimed wave equation, gives \[\psi_0=\frac{x^2}{2s^2}+\frac1s,\quad
\Box_g\psi_0=\frac1{2s^2},\quad
Q_g(\log(\rho/c),\psi_0)=-\frac1{2s^2}.\] Thus paper (2.2) predicts \(1/(4s^2)\), whereas the actual value is \(1/(2s^2)\). Its left-minus-right residual is \(1/(4s^2)\ne0\). Equation [correct] (\#correct) has residual zero.\par
For the same solution, \[w=s^{-1/2}-\frac{x}{2s},\quad
\underline{w}=s^{-1/2}+\frac{x}{2s},\quad \psi_2=0,\] and \[\Box_gw=\Box_g\underline{w}=0,\quad
Q_g(w,w)=Q_g(\underline{w},\underline{w})=0,\quad
Q_g(w,\underline{w})=-\frac1{2s^2}.\] Equation [claimedw] (\#claimedw) predicts \(-1/(8s^{3/2})\), not zero. Its residual is \(1/(8s^{3/2})\). Equation [correctw] (\#correctw) instead gives zero.\par
The example is not a small global Sobolev perturbation of the paper’s rarefaction data: the affine velocity on \(\mathbb R\) is deliberately used only to test a local identity. No assertion is made that it satisfies \((I_0),(I_2),(I_\infty)\), or that it contradicts the main a priori estimate. This distinction is essential to the supported conclusion.\par
\subsection*{Alternative readings and counterevidence}
The manifold has two spatial dimensions, as explicitly stated and as visible in its metric. The three-dimensional-space conformal convention cannot be substituted silently. The line above (3.8), p. 34, uses \(d\log\Omega=(3-\gamma)(dw+d\underline{w})/(2c)\), which agrees with \(\Omega=\rho/c\), not its square; hence redefining \(\Omega\) at p. 23 alone would not reconcile the source.\par
On constant-\(\Omega\) backgrounds the erroneous term vanishes. At \(\gamma=3\) it also vanishes, but that endpoint is outside the stated interval and does not address the example with \(\gamma=2\). These observations explain why some simple checks could miss the error.\par
A material mitigating fact is that the corrected system consists of the same family of quadratic null forms. The structural assertion of Remark 3.1, p. 34, survives: no new derivative order is introduced. The inspected beginnings of Sections 6 and 8 estimate general linear combinations of these forms, with constants absorbed in inequalities. This supports the possibility of a repair of the proof strategy; it does not itself check every commutator, sign, bootstrap constant, or boundary argument.\par
\subsection*{Dependence and supported impact}
Section 6 explicitly starts from (3.7) on p. 55 to prove the lowest-order energy estimates. Section 8 starts from the same equation on p. 65 and constructs differentiated source terms. Section 8.5, pp. 89–90, combines these bounds to close the energy bootstrap used for the main theorem on p. 37. This is a demonstrated core dependency, not a peripheral typographical use of an isolated formula.\par
The issue is therefore classified Major: invalid core derivation in the reviewed author manuscript. It does not establish that the final theorem is false, that all energy bounds fail, or that the journal final version contains the same formulas. Completing the repair requires substituting the corrected sources throughout the affected estimates and checking the full bootstrap closure, including Section 9. No empirical experiment is needed to establish the local algebraic discrepancy; the independent symbolic computation corroborates, rather than replaces, the analytic proof.\par
\section*{A separate check with no issue found}
Lemma 5.2, pp. 53–55, assumes nonnegative smooth \(E,F\) with the stated monotonicities and \[E(t,u)+F(t,u)\leq At^2+B\int_0^uF(t,v)\,dv
+C\int_\delta^t E(\tau,u)\,\frac{d\tau}{\tau}.\] Its conclusion under \(e^{Bu^*}C\leq1\) is \(E+F\leq3Ae^{Bu}t^2\). The Volterra resolvent in \(u\), followed by monotonicity of \(E\) in \(u\), gives \[E+F\leq Ae^{Bu}t^2+D Y(t),\qquad
D=e^{Bu}C,\quad Y(t)=\int_\delta^t E(\tau,u)\,\frac{d\tau}{\tau}.\] Consequently \((Y/t^D)'\leq Ae^{Bu}t^{1-D}\), \(Y(\delta)=0\). Since \(0<D\leq1\), \[Y(t)\leq\frac{Ae^{Bu}}{2-D}
  (t^2-\delta^{2-D}t^D)
\leq\frac{Ae^{Bu}}{2-D}t^2 .\] This even implies \(E+F\leq2Ae^{Bu}t^2\); the paper’s looser intermediate factor two yields its claimed factor three. The auxiliary function \(H\) need not be nonnegative for the integrated Volterra comparison used here. The displayed choice \(a_0=1/(2C_0)\) and \(u_0^*=(\log2)/(2C_0^2)\) on p. 90 satisfies the lemma’s smallness condition. This check does not independently establish all inequalities fed into that application.\par
\section*{Version checks and conclusion}
On the consultation date the arXiv abstract/history page listed v2 of 19 September 2024 as its last revision. The consulted Annals landing page listed the article in volume 202(2), pp. 631–752, and online publication on 26 September 2025. No correction notice was observed on those consulted pages or located by the stated targeted search. This is a limited observation, not proof that no erratum exists elsewhere. The publisher’s final PDF was not compared, so the finding is attributed only to the pinned author manuscript.\par
Issue judgment | Major for (2.2)/(3.7) and their documented core use.\par
Repair verification | Corrected local identities analytically derived and checked on an exact solution. Full repaired theorem proof remains unchecked.\par
Other checked claim | No issue found in Lemma 5.2 within the inspected scope.\par
Adoption | Not requested; no author contact.\par
Scientific acceptance | Pending independent review; no human approval asserted.\par
Remaining work | Unread proof portions, full repair propagation, and journal-version comparison. Report production records are supplied with the targeted continuation.\par
9 T.-W. Luo and P. Yu, On the stability of multi-dimensional rarefaction waves I: the energy estimates, arXiv:2302.09714v2, 19 September 2024. https://arxiv.org/abs/2302.09714v2. Reviewed files are those bound by the local source manifest. Annals of Mathematics, article landing page, DOI 10.4007/annals.2025.202.2.2, https://annals.math.princeton.edu/2025/202-2/p02. Metadata consulted 25 October 2026; final article text not compared.\par
\section*{Supporting checks}
The reviewer-authored script\par
> 00000TOKEN\par
constructs the metric matrix and its inverse from [example] (\#example), differentiates the divergence-form wave operator, and simplifies the residuals with SymPy 1.13.3. It was executed with `python -B` and completed with exit code zero (recorded process 52200). Its machine-readable output is `linearization\_check\_result.json`. All continuity, momentum, potential, enthalpy, and corrected-wave residual assertions passed. No author program was imported or executed.\par
Input validation, extracted page text, rendered source-page images, visual inspection records, source-line fragments, and the exact report hash are retained in the per-paper evidence and provenance files. Text extraction alone is not recorded as visual inspection. Pages 29 and 30 were rendered during investigation but are not counted as viewed.\par
The authorized followup uses the existing Tectonic 0.17.0 runtime. Exact-source build logs, input and PDF hashes, and report-page inspection records are retained in the accompanying provenance and followup evidence. No compiler was installed by this conversation. The source-page images were rendered by Poppler and inspected as recorded.\par
> 00002TOKEN\par
No independent backend-route equivalence or mismatch has been established, and no human mathematical verification or adoption is asserted.\par
\end{document}
